\scanpage{12}

est bilinéaire en $x,y$. \deleted{\unsure{S}}

Comme on a, par transposition de $i:\Delta\longrightarrow\underline{\mathrm{Quad}}(E)$

\[\underline{\mathrm{Quad}}(E)^\vee\xrightarrow{i^\vee}\Delta^\vee,\]

on déduit une forme quadratique composée

\diagram{\begin{tikzcd}[ampersand replacement=\&]
\&\underline{\mathrm{Quad}}(E)^\vee\arrow[dr,"i^\vee"]\&\\
Q_i\ \text{ou }Q_\Sigma:\quad E\arrow[ur,"Q"]\arrow[rr]\&\&\Delta^\vee.
\end{tikzcd}}

Dans le cas d'une structure de similitudes fidèle \inserted{$\Sigma$}, $(\Delta,i)$ ne dépend, à iso. \inserted{can.} unique près, que de $\Sigma$, et il en est donc aussi de $=$ de $Q_i$, ce qui justifie alors la notation $Q_\Sigma$. \deleted{Quand on}

\uline{NB} Quand on se donne une base $e$ de $\Delta$, d'où $\Delta\simeq\Delta^\vee\simeq\underline{\mathcal O}$, alors $Q_\Sigma$ s'identifie à une forme quad. ordinaire sur $E$ i.e. à val. dans $\underline{\mathcal O}$, et ce n'est autre que \deleted{$\unsure{f(x)}$} \inserted{$f=i(e)$.}

% margin note written sideways
\margin{1.5. La donnée d'une structure de similitudes \inserted{(\uline{strictement})} \uline{fidèle} \inserted{sur $E$} \unsure{équivaut|revient} à celle d'un Module inv. $\Delta$ sur $X$, et d'une forme quadratique $Q:E\longrightarrow\Delta^\vee$ telle que \uline{$Q(E)$ engendre} $\Delta^\vee$, i.e. \deleted{\unsure{$Q_x$ est non|$Q$}} sur aucune fibre résiduelle, \deleted{\unsure{est nul}} identiquement nulle.}

2.1. On suppose à partir de maintenant que $E$ est loc. libre de rang 2. Soit

\[L=\underline{\det}E=\bigwedge^2 E\]

son Module déterminantiel, qui est un Module inversible. Notons l'iso. can.

\[E\xrightarrow{\sim} E^\vee\otimes L\]

donné par

\[(x\longmapsto(y\longmapsto x\wedge y)),\quad\text{i.e.}\]

associé à l'accouplement canonique

\[E\times E\longrightarrow L=\bigwedge^2 E:(x,y)\longmapsto x\wedge y.\]
% archive sheet number: 11
